\chapter{Decomposizione spettrale: spezzare per regnare}
\label{ch:metodo_spettrale}

\epigraph{``The eigenvectors of the Laplacian know more about
the graph than you might think.''}{Fan R. K. Chung,
\emph{Spectral Graph Theory}~\parencite{chung1997}}

\lettrine[lines=3]{I}{mmagina} di dover organizzare l'orario
di un istituto comprensivo enorme: 90 classi, 200 docenti,
12 indirizzi diversi. Sembra un'impresa formidabile da
risolvere ``tutta insieme''. Ma se osservi il grafo che ha
per nodi le classi e i docenti, e archi che dicono ``questo
docente insegna in questa classe'', noti qualcosa: il grafo
non \`e omogeneo. Ci sono \emph{cluster} naturali. Le classi
del liceo classico hanno docenti che lavorano quasi
esclusivamente con il classico; le classi dell'ITIS hanno i
loro. I docenti di matematica si distribuiscono fra
indirizzi affini ma raramente toccano l'IPSIA accanto.

Se tu potessi \emph{vedere} questi cluster e risolvere
l'orario di ognuno separatamente, il problema cambierebbe
faccia: invece di una scuola da 90 classi, avresti 5--7
sotto-scuole da 12--18 classi. Ognuna trattabile in pochi
minuti. Questo \`e il trucco della decomposizione spettrale:
\emph{usare l'algebra del grafo per individuare i cluster
naturali}.

\section{Da dove viene l'idea}

\begin{storiabox}
Nel 1973 il matematico ceco Miroslav
Fiedler~\parencite{fiedler1973} pubblic\`o un risultato che
all'epoca sembrava puramente teorico: il \emph{secondo}
autovettore del Laplaciano di un grafo (oggi chiamato
\emph{vettore di Fiedler}\index{Fiedler!vettore}) ne rivela
la struttura comunitaria. Quasi vent'anni dopo Shi e
Malik~\parencite{shi2000} riadattarono l'idea per la
\emph{segmentazione di immagini}: vedi un'immagine come
grafo (pixel = nodi, archi = somiglianza fra pixel
adiacenti), calcoli il vettore di Fiedler e
ottieni i contorni degli oggetti. Da l\`i in poi
il \emph{spectral clustering}~\parencite{vonluxburg2007} \`e
diventato uno standard nella community di machine learning.
\end{storiabox}

L'idea nascosta nel risultato di Fiedler \`e che la
\emph{struttura combinatoria} di un grafo (chi \`e connesso
con chi) si riflette nello \emph{spettro} (gli autovalori)
di una matrice associata. E lo spettro \`e calcolabile in
tempo polinomiale standard.

\section{Costruzione passo-passo}

\subsection{(1) Il grafo bipartito classe--docente}

\begin{defbox}
Per una scuola con $n$ classi e $m$ docenti, costruiamo un
\emph{grafo bipartito}\index{grafo bipartito} $G = (V, E)$
con:
\begin{itemize}
\item $V = V_{\text{C}} \cup V_{\text{D}}$, dove
  $V_{\text{C}}$ \`e l'insieme delle classi e $V_{\text{D}}$
  l'insieme dei docenti.
\item $E$: c'\`e un arco $\{c, d\}$ con peso $w_{cd}$ se il
  docente $d$ insegna $w_{cd}$ ore alla classe $c$.
\end{itemize}
\end{defbox}

Esempio in miniatura: una scuola con 3 classi ($c_1, c_2, c_3$)
e 4 docenti ($d_1, d_2, d_3, d_4$), in cui $d_1, d_2$
insegnano solo nelle prime due classi e $d_3, d_4$ solo
nella terza, ha due cluster evidenti: $\{c_1, c_2, d_1, d_2\}$
e $\{c_3, d_3, d_4\}$. Il grafo bipartito ha pochissimi (o
zero) archi fra i due cluster.

\subsection{(2) La matrice di adiacenza e quella dei gradi}

\begin{defbox}
La \emph{matrice di adiacenza} $A$ \`e una matrice quadrata
$|V| \times |V|$ con $A_{ij} = w_{ij}$ se esiste l'arco
$\{i, j\}$, 0 altrimenti.

La \emph{matrice dei gradi} $D$ \`e diagonale con
$D_{ii} = \sum_j A_{ij}$, ovvero la somma dei pesi degli
archi incidenti su $i$. Per un docente, $D_{ii}$ \`e il monte
ore totale; per una classe, \`e la somma delle ore di lezione
ricevute.
\end{defbox}

\subsection{(3) Il Laplaciano (normalizzato)}

\begin{defbox}
Il \emph{Laplaciano}\index{Laplaciano!normalizzato}
\emph{normalizzato} di $G$ \`e:
\[
  L = I - D^{-1/2} A\, D^{-1/2}
\]
dove $I$ \`e la matrice identit\`a.
\end{defbox}

\textbf{Come si legge}: prendi la matrice di adiacenza $A$
(che dice chi \`e connesso a chi), la \emph{normalizzi}
dividendo ogni elemento per la radice del prodotto dei gradi
dei due estremi (per togliere l'effetto della diversa
``importanza'' dei nodi), e sottrai il risultato dall'identit\`a.

\textbf{Propriet\`a chiave}:
\begin{itemize}
\item $L$ \`e simmetrica e semi-definita positiva.
\item Tutti i suoi autovalori $\lambda_0, \lambda_1, \ldots$
  sono nell'intervallo $[0, 2]$.
\item L'autovalore minimo $\lambda_0$ vale sempre 0; il
  corrispondente autovettore \`e proporzionale a
  $D^{1/2}\mathbf{1}$ (banale).
\item Il \emph{secondo autovalore} $\lambda_1$ \`e detto
  \emph{algebraic connectivity}~\parencite{fiedler1973}: \`e
  ``piccolo'' quando il grafo ha cluster ben separati,
  ``grande'' quando il grafo \`e ben connesso.
\end{itemize}

\subsection{(4) Il vettore di Fiedler}

\begin{defbox}
Il \emph{vettore di Fiedler} $v_1$ \`e l'autovettore
associato al secondo autovalore $\lambda_1$ di $L$.
\end{defbox}

L'idea cruciale: ogni nodo del grafo corrisponde a una
componente di $v_1$, cio\`e a un numero reale. Se ordini i
nodi per il valore della loro componente in $v_1$, ottieni
un \emph{ordinamento naturale} che separa le comunit\`a.

\sidenote{Esempio:
se $v_1 = (-0.4, -0.5, +0.3, +0.6)^\top$, l'ordinamento per
componenti riordina i nodi come $\{2, 1, 3, 4\}$. Tagliando
fra il 2\textsuperscript{o} e il 3\textsuperscript{o} ottieni
i due cluster.}

\subsection{(5) k-means sui top-$k$ autovettori}

Per ottenere $k$ cluster (non solo 2), si calcolano i primi
$k$ autovettori non-banali $v_1, v_2, \ldots, v_k$ di $L$.
Si ottiene cos\`i una matrice $|V| \times k$: ogni nodo
diventa un punto in $\mathbb{R}^k$. A questo punto si applica
$k$-means standard su questa rappresentazione \emph{spettrale}
e si ottengono $k$ cluster.

\begin{figure}[h]
\centering
\begin{tikzpicture}[scale=0.9, every node/.style={font=\small\sffamily}]
  % cluster A
  \node[circle, draw, fill=cPref!30, minimum size=8mm] (c1) at (0,2) {$c_1$};
  \node[circle, draw, fill=cPref!30, minimum size=8mm] (c2) at (0,1) {$c_2$};
  \node[circle, draw, fill=cPref!10, minimum size=8mm] (d1) at (1.5,2) {$d_1$};
  \node[circle, draw, fill=cPref!10, minimum size=8mm] (d2) at (1.5,1) {$d_2$};
  \draw (c1)--(d1); \draw (c1)--(d2); \draw (c2)--(d1); \draw (c2)--(d2);
  % cluster B
  \node[circle, draw, fill=cOro!30, minimum size=8mm] (c3) at (5,2) {$c_3$};
  \node[circle, draw, fill=cOro!30, minimum size=8mm] (c4) at (5,1) {$c_4$};
  \node[circle, draw, fill=cOro!10, minimum size=8mm] (d3) at (6.5,2) {$d_3$};
  \node[circle, draw, fill=cOro!10, minimum size=8mm] (d4) at (6.5,1) {$d_4$};
  \draw (c3)--(d3); \draw (c3)--(d4); \draw (c4)--(d3); \draw (c4)--(d4);
  % bridge
  \draw[dashed, thin] (d2)--(c3) node[midway, above, font=\scriptsize, text=cSiena] {ponte};
  % rectangles
  \draw[thick, rounded corners=4pt, cPref] ($(c1.north west)+(-0.3,0.3)$) rectangle ($(d2.south east)+(0.3,-0.3)$);
  \draw[thick, rounded corners=4pt, cOro] ($(c3.north west)+(-0.3,0.3)$) rectangle ($(d4.south east)+(0.3,-0.3)$);
  \node[font=\sffamily\bfseries, color=cPref] at (0.75, 3) {Cluster A};
  \node[font=\sffamily\bfseries, color=cOro] at (5.75, 3) {Cluster B};
\end{tikzpicture}
\caption{Esempio di grafo bipartito con due cluster naturali
e un ``ponte'' debole. Il vettore di Fiedler assegna
componenti di segno opposto ai due cluster.}
\label{fig:spettrale-cluster}
\end{figure}

\section{Esempio mini-completo a mano}

Prendiamo la scuola di fig.~\ref{fig:spettrale-cluster}:
4 classi, 4 docenti, archi come disegnati (incluso un
``ponte'' debole $d_2$--$c_3$).

\textbf{Passo 1: matrici.} Numeriamo i nodi
$1, 2, \ldots, 8$ in ordine $c_1, c_2, c_3, c_4, d_1, d_2,
d_3, d_4$. La matrice di adiacenza $A$ \`e (con peso 1 ovunque):
\[
A = \begin{pmatrix}
0 & 0 & 0 & 0 & 1 & 1 & 0 & 0 \\
0 & 0 & 0 & 0 & 1 & 1 & 0 & 0 \\
0 & 0 & 0 & 0 & 0 & 1 & 1 & 1 \\
0 & 0 & 0 & 0 & 0 & 0 & 1 & 1 \\
1 & 1 & 0 & 0 & 0 & 0 & 0 & 0 \\
1 & 1 & 1 & 0 & 0 & 0 & 0 & 0 \\
0 & 0 & 1 & 1 & 0 & 0 & 0 & 0 \\
0 & 0 & 1 & 1 & 0 & 0 & 0 & 0 \\
\end{pmatrix}
\]

I gradi sono: $D = \mathrm{diag}(2, 2, 3, 2, 2, 3, 2, 2)$.

\textbf{Passo 2: Laplaciano normalizzato.} Si ottiene
applicando la formula $L = I - D^{-1/2} A D^{-1/2}$ (il
calcolo a mano \`e tedioso, lo demandiamo a NumPy).

\textbf{Passo 3: spettro.} Calcolando con
\texttt{numpy.linalg.eigh(L)} si ottiene:
\[
  \lambda_0 \approx 0,\quad
  \lambda_1 \approx 0.18,\quad
  \lambda_2 \approx 0.72,\quad \ldots
\]
Il piccolo valore $\lambda_1 \approx 0.18$ conferma che il
grafo \`e debolmente connesso: i due cluster sono ben
separati, il ponte $d_2$--$c_3$ \`e l'unico legame.

\textbf{Passo 4: vettore di Fiedler.} \`E l'autovettore
associato a $\lambda_1$. Calcolandolo si trova qualcosa come:
\[
  v_1 \approx (+0.40, +0.40, -0.35, -0.40, +0.35, +0.30, -0.30, -0.30)^\top
\]
I primi quattro nodi (cluster A: $c_1, c_2, d_1, d_2$) hanno
componenti positive; gli altri quattro (cluster B:
$c_3, c_4, d_3, d_4$) hanno componenti negative. La
separazione \`e netta. Tagliando in zero, otteniamo
esattamente i due cluster.

\begin{esempiobox}
Per ottenere $k > 2$ cluster, si usano i primi $k$
autovettori e si applica $k$-means. Su questa stessa scuola
con $k = 2$, $k$-means su $v_1$ produce gli stessi due
cluster ottenuti dal segno.
\end{esempiobox}

\section{Generalizzazione: cosa cambia su scuola reale}

In una scuola con 90 classi e 200 docenti il calcolo \`e
identico, solo che le matrici sono $290 \times 290$ e gli
autovettori si calcolano in qualche secondo con
LAPACK/scipy. La differenza pratica \`e:

\begin{itemize}
\item I cluster non sono \emph{cos\`i} netti come nell'esempio
  giocattolo. Tipicamente ci sono cluster ben separati e
  poi ``nodi di confine'' che potrebbero stare in pi\`u
  cluster.
\item Il numero $k$ di cluster non \`e dato a priori: si
  sceglie con una \emph{eigengap heuristic}: si cerca il
  $k$ tale che $\lambda_k$ sia molto pi\`u grande di
  $\lambda_{k-1}$ (un ``salto'' nello spettro). $\pi$Tantum
  esegue un'analisi automatica del gap e suggerisce un
  valore di $k$, ma l'utente pu\`o forzarlo.
\item Una volta definiti i cluster, la \emph{ricucitura}
  ($\pi$Tantum la chiama \emph{Phase~B stitch}) richiede un
  passo CP-SAT globale che armonizza i ``ponti'' fra
  cluster (lezioni che coinvolgono nodi di cluster diversi).
\end{itemize}

\section{Quando funziona, quando no}

\begin{ideabox}
La decomposizione spettrale funziona benissimo quando
\emph{i cluster naturali esistono}: scuole con indirizzi
distinti, classi-di-concorso disgiunte, presenza di sedi
distaccate. Il guadagno tipico su istanze grandi \`e
$5\times$--$10\times$ rispetto a CP-SAT monolitico.
\end{ideabox}

\begin{avvertenzabox}
Quando il grafo \`e \emph{denso} (i docenti insegnano un po'
ovunque, niente cluster naturali), il vettore di Fiedler
non separa nulla e i cluster prodotti sono artificiali. In
quel caso la decomposizione genera ``ponti'' molti e pesanti,
e la fase di ricucitura diventa pi\`u costosa della
risoluzione monolitica. $\pi$Tantum lo rileva guardando
$\lambda_1$ del Laplaciano: se \`e troppo grande (sopra
$\sim 0.5$), suggerisce di saltare la decomposizione.
\end{avvertenzabox}

\subsection{Parametri principali}

\begin{itemize}
\item \texttt{n\_clusters}: numero di cluster da produrre.
  Default: scelto via eigengap heuristic; tipicamente
  4--8 per scuole medio-grandi.
\item \texttt{laplacian\_type}: \texttt{normalized}
  (default) o \texttt{unnormalized}. Il primo \`e
  raccomandato per grafi non bilanciati nei pesi.
\item \texttt{kmeans\_init}: \texttt{k-means++} (default,
  robusto) o \texttt{random}.
\item \texttt{stitch\_passes}: numero di iterazioni della
  fase di ricucitura. Default: 1.
\end{itemize}

\section{Connessioni}

\begin{connessionebox}
\begin{itemize}
\item \textbf{CP-SAT} (cap.~\ref{ch:metodo_cpsat}): ogni
  cluster \`e risolto da CP-SAT come problema autonomo.
\item \textbf{Hall pre-check} (cap.~\ref{ch:metodo_hall}): si
  esegue \emph{su ciascun cluster} oltre che globalmente; se
  un cluster non passa Hall, si rilassano i ponti per
  ridistribuire la domanda.
\item \textbf{Lagrangian Relaxation}
  (cap.~\ref{ch:metodo_lagrangian}): un'alternativa alla
  ricucitura combinatoria \`e la dualizzazione dei vincoli
  inter-cluster con moltiplicatori di Lagrange.
\item \textbf{Analisi del grafo} (cap.~\ref{ch:metodo_grafo}):
  modularit\`a, betweenness, densit\`a sono \emph{misure} che
  $\pi$Tantum espone in dashboard per dare al coordinatore
  un'idea della ``segmentabilit\`a'' della sua scuola
  \emph{prima} di lanciare la pipeline.
\end{itemize}
\end{connessionebox}

\begin{biblioparvabox}
\begin{itemize}
\item \cite{fiedler1973}: il paper originale di Fiedler.
  Tre pagine, ancora leggibili.
\item \cite{chung1997}: il libro di riferimento sulla teoria
  spettrale dei grafi. Tecnico ma chiaro.
\item \cite{vonluxburg2007}: tutorial divulgativo sullo
  spectral clustering, con tutti i dettagli pratici (scelta
  di $k$, normalizzazioni, $k$-means vs alternative).
\item \cite{shi2000}: l'applicazione al \emph{normalized cut}
  per immagini, che ha reso popolare l'idea fuori dalla
  matematica pura.
\end{itemize}
\end{biblioparvabox}
